%%%%%%%%%%%%%%%%%%%%%%%%%%%%%%%%%%%%%%%%%%%%%%%%%%%
%
%  New template code for TAMU Theses and Dissertations starting Spring 2021.
%
%
%  Last Updated: 1/13/2021
%
%%%%%%%%%%%%%%%%%%%%%%%%%%%%%%%%%%%%%%%%%%%%%%%%%%%
%%%%%%%%%%%%%%%%%%%%%%%%%%%%%%%%%%%%%%%%%%%%%%%%%%%%%%%%%%%%%%%%%%%%%
%%                           ABSTRACT
%%%%%%%%%%%%%%%%%%%%%%%%%%%%%%%%%%%%%%%%%%%%%%%%%%%%%%%%%%%%%%%%%%%%%

\chapter*{ABSTRACT}
\addcontentsline{toc}{chapter}{ABSTRACT} % Needs to be set to part, so the TOC doesn't add 'CHAPTER ' prefix in the TOC.

\pagestyle{plain} % No headers, just page numbers
\pagenumbering{roman} % Roman numerals
\setcounter{page}{2}

Weaving is one of humanity's oldest activities, involving the creation of structures by interlacing threads or filaments. It exemplifies a fundamental principle: the whole is greater than the sum of its parts. Woven objects display material properties and functions not present in the individual threads. Applications of weaving range from everyday objects such as baskets and fabrics to advanced fields such as construction and high-performance composites. Moreover, similar structural principles appear in nature, where woven-like arrangements are observed in elastic biological tissues.

Woven structures consist of interlaced threads forming ordered, repeating patterns. Beyond textiles, they are valued for their structural efficiency and appear in architecture, engineered composites, and digital fabrication. Traditional fabrics are anisotropic, exhibiting strength primarily along thread directions in a plane. Extending textiles to industrial and engineering applications---where strength is needed in multiple directions---led to the development of three-dimensional woven textiles. These volumetric weaves, which generalize fabrics into three dimensions, offer new opportunities for material design but currently lack a unified theoretical framework. From a manufacturing perspective, the simplest approach has been to build these textiles layer by layer, connecting layers with straight threads. Despite advances in manufacturing, there remains a gap in systematically designing and understanding three-dimensional woven structures. This dissertation addresses that gap, starting from a single thread and connecting concepts to develop a framework for designing volumetric woven textiles.

To study three-dimensional weaving, we begin by revisiting planar weaves. In mathematics, planar weaves can be represented using square or triangular grids depending on the pattern. To better understand these structures, we modularize the system by tiling the 2D space. The first part of this dissertation addresses the problem of space-filling threads. Previous approaches produced fixed, non-reconfigurable threads, whereas our modular tiles allow reconfiguration through translation and rotation, generating infinite families of bidirectional planar weaves. Beyond reducing the need for filler elements, this modularity highlights the role of symmetry in fabric design, providing a foundation for extending these principles into volumetric weaving.

Second, we develop an algebraic framework for nonwoven volumetric textiles. Nonwoven structures differ from woven ones in that they do not strictly follow the over--under interlacing rule; instead, they are networks of paths that still achieve fabric-like connectivity. By extending these ideas into three dimensions, we show how we can generate symbols such that we can identify families of structures by rules rather than by direct construction, a process we call algebraic parametrization. This approach allows the creation of configurations that mimic the mechanical behavior of woven fabrics while also realizing structures not possible in planar weaves.

Finally, we formalize volumetric weaves using the isometries of Bravais lattices. Extending the modularity concept to three dimensions allows us to define volumetric weaves as assemblies of repeatable modules, analogous to the regularity observed in crystals. This lattice-based framework provides a systematic design space for volumetric weaves, enabling locally configurable structures that go beyond current understanding and achieve complexity far exceeding that of planar weaves.

Together, these studies establish a framework for symmetry-driven, modular design of volumetric weaves. By integrating geometric invariance (properties unchanged under symmetry transformations), algebraic parametrization (symbolic generation of structure families), and lattice isometries (crystalline symmetries as organizing principles), this work provides a unified foundation for three-dimensional weaving. The framework has applications in engineered materials, composite structures, and digitally fabricated architectures. More broadly, it demonstrates how weaving, traditionally a planar craft, can be extended into a systematic science of volumetric interlacement, expanding possibilities for both functional and aesthetic material design.

\pagebreak{}
